\section{Match coefficient order to window order}
Let \(x_0\) be the first new sample. The supplied array begins with
\(x_{-7},\ldots,x_{-1}\), then \(x_0,x_1,\ldots\). Its first output uses all
positions zero through seven. The next output starts one sample later and uses
positions one through eight. Advance the full input address by two bytes for each
output, not by sixteen bytes.

Some descriptions write the filter as \(y_i=\sum_{k=0}^{7}h_kx_{i-k}\), where
\(h_0\) multiplies the current sample. In this interface \(h_k=c_{7-k}\).
The stored coefficient array is reversed relative to that conventional current-first
notation, so the two loaded windows align in memory order. Reversing both arrays is
unnecessary; interpreting the stored order correctly is essential.

\begin{figure}[htbp]
\centering
\begin{tikzpicture}[alt={Two chronological windows share seven samples. Output zero uses array positions zero through seven, with coefficients c zero through c seven. Output one uses positions one through eight with the same coefficient sequence. Array position seven is the first new sample.},x=0.82cm,y=0.7cm,font=\small]
\foreach \j in {0,...,8}{
 \draw (\j,0) rectangle ++(1,0.8);
 \node at (\j+0.5,0.4) {\j};
}
\draw[thick,MARSAccent] (0,1.05)--(8,1.05);
\draw[thick,MARSAccent] (1,-0.25)--(9,-0.25);
\node[anchor=east] at (-0.15,1.05) {Output 0};
\node[anchor=east] at (0.85,-0.65) {Output 1};
\node[anchor=west] at (0,1.7) {Seven history samples: positions 0--6};
\draw[-{Stealth}] (7.5,1.55)--(7.5,0.9);
\node[anchor=west] at (8,1.7) {Current};
\end{tikzpicture}
\caption{The first two eight-tap input windows. Coefficients keep the same order as the window advances one sample.}
\end{figure}

The final output begins at input position \(n-1\) and ends at position \(n+6\).
Thus \(n+7\) input halfwords suffice, including the history. A one-output call still
has a complete sixteen-byte input window. The previous chapter's eight-output block
gate does not apply: this vectorization runs across taps within one output.
